\documentclass[12pt,a4paper]{article}

\usepackage[utf8]{inputenc}
\usepackage[T1]{fontenc}
\usepackage[spanish,es-tabla]{babel}
\usepackage{geometry}
\usepackage{hyperref}
\usepackage{lmodern}
\usepackage{setspace}
\usepackage{fancyhdr}
\usepackage{longtable}
\usepackage{array}
\usepackage{enumitem}

\geometry{margin=2.5cm}
\setstretch{1.15}
\setlength{\headheight}{15pt}
\pagestyle{fancy}
\fancyhf{}
\lhead{Proyecto DAW}
\rhead{App de reservas de pistas}
\cfoot{\thepage}

\title{Desarrollo de una Aplicación Web para la Gestión de Reservas de Pistas}\author{Alumno: \underline{Tu nombre aquí} \\
Ciclo Formativo de Grado Superior en Desarrollo de Aplicaciones Web}
\date{\today}

\begin{document}

% -------------------- Portada --------------------
\pagenumbering{gobble}
\begin{titlepage}
    \centering
    \thispagestyle{empty}
    \vspace*{2cm}

    {\Large\textbf{App Final}}\\[1.2cm]
    {\huge\textbf{Desarrollo de una Aplicación Web para la Gestión de Reservas de Pistas}}\\[1.5cm]

    \textbf{Asignaturas implicadas:}\\[0.3cm]
    Cliente Servidor\\
    Contenedores\\
    Desarrollo de Aplicaciones Web\\[2cm]

    \begin{tabular}{>{\bfseries}p{5cm} p{7cm}}
        Alumno & Sebastià Romaguera \\
        Centro & Borja Moll \\
        Profesor & Alejandro Alberti De La Rosa \\
        Curso & 2 DAW Intensivo \\
        Fecha & Mayo 2026 \\
    \end{tabular}

    \vfill
\end{titlepage}

\clearpage
\pagenumbering{arabic}
\tableofcontents
\newpage

% -------------------- Introduccion --------------------
\section{Introducción y Alcance}

La idea de este proyecto es desarrollar una aplicación web inspirada en plataformas de reserva deportiva como Playtomic, pero adaptada a un alcance más sencillo y centrado en la gestión de pistas y partidos, sin incluir pagos ni funcionalidades comerciales avanzadas.

El problema que se quiere resolver tiene dos partes. Por un lado, plataformas como Playtomic aplican costes de gestión por reserva que encarecen la experiencia del usuario. Este proyecto busca ofrecer una alternativa más económica, donde la gestión de reservas sea gratuita. Por otro lado, se mantiene claro que el uso de la pista en cada club sí se paga, ya que ese coste depende de la instalación deportiva y no de la plataforma.

Además, en muchos clubes o instalaciones deportivas la gestión sigue haciéndose de forma manual, por teléfono, por mensajería o con sistemas poco intuitivos. Esto provoca errores, duplicidades, pérdida de tiempo y una mala experiencia tanto para los usuarios como para los responsables de la instalación.

La aplicación propuesta permitirá que los usuarios puedan consultar pistas disponibles, reservar una franja horaria, crear o unirse a partidos y revisar su actividad. Desde el punto de vista técnico, el proyecto servirá para integrar una arquitectura cliente-servidor, una base de datos relacional y un despliegue mediante contenedores, por lo que encaja de forma directa con varias de las competencias del ciclo.

\subsection{Objetivos del proyecto}
\begin{itemize}[leftmargin=*]
    \item Centralizar la reserva de pistas en una sola plataforma web.
    \item Facilitar la consulta de disponibilidad en tiempo real.
    \item Reducir la gestión manual de reservas y conflictos de horario.
    \item Permitir la organización básica de partidos entre usuarios.
    \item Aplicar una arquitectura moderna con tecnologías actuales de desarrollo web.
    \item Desplegar la solución de forma reproducible mediante contenedores.
\end{itemize}

\subsection{Alcance funcional}
Para mantener el proyecto realista dentro del tiempo de desarrollo de un trabajo de FP, el alcance se limita a las siguientes funciones principales:
\begin{itemize}[leftmargin=*]
    \item Registro e inicio de sesión de usuarios.
    \item Listado de pistas y disponibilidad.
    \item Reserva de pistas por franjas horarias.
    \item Creación y gestión básica de partidos.
    \item Consulta del historial de reservas.
\end{itemize}

Quedan fuera del alcance, al menos en esta primera versión, funciones como pagos integrados, chat en tiempo real, recomendación automática de partidos, sistema de valoraciones o integraciones externas complejas.

% -------------------- Pila de tecnologias --------------------
\section{Pila de tecnologías}

La selección de tecnologías se ha realizado buscando equilibrio entre facilidad de desarrollo, claridad para la memoria técnica y adecuación al contenido de las asignaturas implicadas.

\subsection{Frontend}
\begin{itemize}[leftmargin=*]
    \item \textbf{React + TypeScript}: permite construir una interfaz modular, mantenible y tipada.
    \item \textbf{Vite}: proporciona un entorno de desarrollo rápido y un arranque ligero.
    \item \textbf{CSS moderno}: para el diseño visual y la adaptación a distintos tamaños de pantalla.
\end{itemize}

El frontend será el responsable de mostrar la información al usuario, gestionar formularios, navegar entre vistas y consumir la API del backend.

\subsection{Backend}
\begin{itemize}[leftmargin=*]
    \item \textbf{Node.js}: entorno de ejecución en JavaScript para el servidor.
    \item \textbf{TypeScript}: aporta tipado estático y mejora la mantenibilidad del código.
    \item \textbf{Express}: framework ligero para crear la API REST de manera sencilla.
\end{itemize}

El backend actuará como intermediario entre la interfaz de usuario y la base de datos. Su función principal será exponer endpoints para autenticación, reservas, partidos y consulta de datos.

\subsection{Base de datos}
\begin{itemize}[leftmargin=*]
    \item \textbf{PostgreSQL}: base de datos relacional robusta, muy adecuada para gestionar usuarios, pistas, reservas y relaciones entre entidades.
\end{itemize}

\subsection{Contenedores}
\begin{itemize}[leftmargin=*]
    \item \textbf{Docker}: permitirá empaquetar frontend, backend y base de datos en entornos aislados.
    \item \textbf{Docker Compose}: facilitará el arranque conjunto de todos los servicios con una única configuración.
\end{itemize}

El uso de contenedores es especialmente interesante para la asignatura de Contenedores porque permite reproducir el proyecto en cualquier máquina sin depender de instalaciones locales complejas.

\subsection{Herramientas de apoyo}
\begin{itemize}[leftmargin=*]
    \item \textbf{Git y GitHub}: control de versiones y alojamiento del código.
    \item \textbf{Postman}: pruebas de la API.
    \item \textbf{VS Code}: entorno principal de desarrollo.
\end{itemize}

\subsection{Resumen de la arquitectura}
La aplicación seguirá una arquitectura cliente-servidor en la que el navegador del usuario consumirá una API REST. El frontend se encargará de la experiencia de usuario, mientras que el backend procesará la lógica de negocio y la persistencia se gestionará en PostgreSQL.

\end{document}
